\documentclass[a4paper]{article}

\usepackage[spanish]{babel}
\usepackage{graphicx}
\usepackage{amsmath, amssymb}
\usepackage[margin=2cm]{geometry}
\usepackage{fancyhdr}
\usepackage{enumerate}
\usepackage[shortlabels]{enumitem}
\usepackage{parskip}
\usepackage[most]{tcolorbox}
\usepackage[hidelinks, colorlinks=true, linkcolor=green!60!black, urlcolor=green!60!black]{hyperref}
\usepackage{float}
\usepackage{booktabs}
\usepackage{mathrsfs}


% cabecera
\pagestyle{fancy}
\fancyhead[l]{Integral}
\fancyhead[c]{Función Beta}
\fancyhead[r]{\today}
\fancyfoot[c]{\thepage}
\renewcommand{\headrulewidth}{0.2pt} % linea horizontal


\begin{document}

\tableofcontents

\section{Integral $I = \int_0^{\pi/2} \ln^2(\cos x)\, dx$ mediante la función Beta}

Evaluar la integral definida:

$$I = \int_0^{\pi/2} \ln^2(\cos x)\, dx$$

utilizando la expresión integral de la función Beta y diferenciación bajo el signo integral.

\subsection{Solución}

Recordemos la expresión integral de la función Beta:

$$B(u,v) = 2\int_0^{\pi/2} \cos^{2u-1}x\, \sin^{2v-1}x\, dx$$

Derivamos parcialmente respecto a $u$:

$$\partial_u B = 2\int_0^{\pi/2} \partial_u\!\left(\cos^{2u-1}x\right) \sin^{2v-1}x\, dx$$

Calculamos la derivada del factor trigonométrico:

$$\partial_u\!\left(\cos^{2u-1}x\right) = \partial_u\!\left(e^{(2u-1)\ln(\cos x)}\right) = 2\ln(\cos x)\cdot \cos^{2u-1}x$$

Por lo tanto:

$$\partial_u B = 4\int_0^{\pi/2} \cos^{2u-1}x\, \sin^{2v-1}x\, \ln(\cos x)\, dx$$

Derivamos nuevamente y obtenemos:

$$\partial_u^2 B = 8\int_0^{\pi/2} \cos^{2u-1}x\, \sin^{2v-1}x\, \ln^2(\cos x)\, dx$$

Con esta expresión podemos construir nuestra integral haciendo los exponentes de las funciones trigonométricas igual a cero:

$$2u - 1 = 0 \implies u = \tfrac{1}{2}$$
$$2v - 1 = 0 \implies v = \tfrac{1}{2}$$

De modo que:

$$I = \int_0^{\pi/2} \ln^2(\cos x)\, dx = \frac{1}{8}\,\partial_u^2 B(u,v)\Big|_{u=v=1/2}$$

Recordemos la relación entre la función Beta y la función Gamma:

$$B(u,v) = \frac{\Gamma(u)\,\Gamma(v)}{\Gamma(u+v)}$$

Derivando respecto a $u$:

$$\partial_u B = \frac{\Gamma'(u)\,\Gamma(v)\,\Gamma(u+v) - \Gamma(u)\,\Gamma(v)\,\Gamma'(u+v)}{\Gamma^2(u+v)}$$

Recordemos la función digamma:

$$\psi(s) = \frac{d}{ds}\big[\log(\Gamma(s))\big] = \frac{\Gamma'(s)}{\Gamma(s)} \implies \Gamma'(s) = \psi(s)\,\Gamma(s)$$

Sustituyendo y factorizando:

$$\partial_u B = \frac{\Gamma(v)\,\Gamma(u)\,\Gamma(u+v)}{\Gamma^2(u+v)}\big[\psi(u) - \psi(u+v)\big]$$

Simplificando:

$$\partial_u B = B(u,v)\big[\psi(u) - \psi(u+v)\big]$$

Notación: $\psi'(s) = \psi_1(s)$ (función trigamma). Derivando nuevamente:

$$\partial_u^2 B = \partial_u B\big[\psi(u) - \psi(u+v)\big] + B\big[\psi_1(u) - \psi_1(u+v)\big]$$

$$\partial_u B\Big|_{u=v=1/2} = B\!\left(\tfrac{1}{2},\tfrac{1}{2}\right)\left[\psi\!\left(\tfrac{1}{2}\right) - \psi(1)\right]$$

Calculamos cada valor:

$$B\!\left(\tfrac{1}{2},\tfrac{1}{2}\right) = \frac{\Gamma(1/2)\,\Gamma(1/2)}{\Gamma(1)} = \frac{\sqrt{\pi}\cdot\sqrt{\pi}}{1} = \pi$$

$$\psi\!\left(\tfrac{1}{2}\right) = \frac{\Gamma'(1/2)}{\Gamma(1/2)} = \frac{-\sqrt{\pi}\,(\gamma + 2\ln 2)}{\sqrt{\pi}} = -(\gamma + 2\ln 2)$$

$$\psi(1) = \frac{\Gamma'(1)}{\Gamma(1)} = -\gamma$$

Sustituyendo:

$$\partial_u B\Big|_{u=v=1/2} = \pi\big[-\gamma - 2\ln 2 + \gamma\big] = -2\pi\ln 2$$

Escribamos $\psi_1$ como serie partiendo de la serie de $\psi$:

$$\psi(1+z) = -\gamma + \sum_{k\geq 1}\left(\frac{1}{k} - \frac{1}{k+z}\right)$$

Derivando:

$$\psi_1(1+z) = \sum_{k\geq 1}\frac{1}{(k+z)^2}$$

Evaluando en $z = -\tfrac{1}{2}$:

$$\psi_1\!\left(\tfrac{1}{2}\right) = \sum_{k\geq 1}\frac{1}{\left(k - \tfrac{1}{2}\right)^2} = 4\sum_{k\geq 1}\frac{1}{(2k-1)^2}$$

Recordemos el problema de Basilea:

$$\frac{\pi^2}{6} = \sum_{k\geq 1}\frac{1}{k^2} = \sum_{k\geq 1}\frac{1}{(2k)^2} + \sum_{k\geq 1}\frac{1}{(2k-1)^2} = \frac{1}{4}\cdot\frac{\pi^2}{6} + \sum_{k\geq 1}\frac{1}{(2k-1)^2}$$

Despejando:

$$\sum_{k\geq 1}\frac{1}{(2k-1)^2} = \frac{\pi^2}{6} - \frac{\pi^2}{24} = \frac{\pi^2}{8}$$

Por lo tanto:

$$\psi_1\!\left(\tfrac{1}{2}\right) = 4\cdot\frac{\pi^2}{8} = \frac{\pi^2}{2}$$


Evaluando la serie en $z = 0$:

$$\psi_1(1) = \sum_{k\geq 1}\frac{1}{k^2} = \frac{\pi^2}{6}$$


Sustituimos todos los valores en la expresión de la segunda derivada:

$$\partial_u^2 B\Big|_{u=v=1/2} = -2\pi\ln 2\left[\psi\!\left(\tfrac{1}{2}\right) - \psi(1)\right] + \pi\left[\psi_1\!\left(\tfrac{1}{2}\right) - \psi_1(1)\right]$$

$$= -2\pi\ln 2\big[-2\ln 2\big] + \pi\left(\frac{\pi^2}{2} - \frac{\pi^2}{6}\right)$$

$$= 4\pi\ln^2 2 + \pi\cdot\frac{\pi^2}{3} = 4\pi\ln^2 2 + \frac{\pi^3}{3}$$


$$I = \frac{1}{8}\,\partial_u^2 B(u,v)\Big|_{u=v=1/2} = \frac{1}{8}\left(4\pi\ln^2 2 + \frac{\pi^3}{3}\right) = \frac{\pi}{2}\ln^2 2 + \frac{\pi^3}{24}$$

$$\boxed{\int_0^{\pi/2} \ln^2(\cos x)\, dx = \frac{\pi}{2}\left(\ln^2 2 + \frac{\pi^2}{12}\right)}$$

\subsection{Comprobación computacional}

Para verificar el resultado obtenido, puede consultar las siguientes referencias en WolframAlpha:

\begin{itemize}
    \item \href{https://www.wolframalpha.com/input?i=Integrate%5BLog%5BCos%5Bx%5D%5D%5E2%2C%7Bx%2C0%2CPi%2F2%7D%5D}{Verificación de la integral $\displaystyle\int_0^{\pi/2} \ln^2(\cos x)\, dx$}
    
    \item \href{https://www.wolframalpha.com/input?i=Pi%2F2*%28Log%5B2%5D%5E2%2B%5Cpi%5E2%2F12%29}{Evaluación del resultado $\displaystyle\frac{\pi}{2}\left(\ln^2 2 + \frac{\pi^2}{12}\right)$}
\end{itemize}

\end{document}